\documentclass[a4paper,11pt]{article}
\usepackage[utf8]{inputenc}
\usepackage[T1]{fontenc}
\usepackage[french]{babel}
\usepackage[left=2cm,right=2cm,top=2cm,bottom=2cm]{geometry}
\usepackage{xcolor}
\usepackage{hyperref}
\usepackage{fontawesome5}
\usepackage{parskip}

% --- Couleurs Wevioo ---
\definecolor{primary}{RGB}{0, 100, 160}
\hypersetup{colorlinks=true, urlcolor=primary}

\begin{document}

% --- EXPÉDITEUR ---
\noindent
\begin{minipage}[t]{0.5\textwidth}
    \textbf{Loïc Aron MBASSI EWOLO}\\
    Élève-Ingénieur Bac+5 -- Double diplôme ENSEA\\[3pt]
    \faMapMarker\ Cergy, France\\
    \faPhone\ (+33) 7 59 52 33 98\\
    \faEnvelope\ \href{mailto:loicaron.mbassiewolo@ensea.fr}{loicaron.mbassiewolo@ensea.fr}
\end{minipage}
\hfill
% --- DATE ---
\begin{minipage}[t]{0.4\textwidth}
    \raggedleft
    Cergy, le 9 mars 2026
\end{minipage}

\vspace{1.2cm}

% --- DESTINATAIRE ---
\hfill
\begin{minipage}[t]{0.5\textwidth}
    \raggedleft
    \textbf{Wevioo}\\
    Service Recrutement\\
    Paris, France
\end{minipage}

\vspace{1cm}

\noindent\textbf{Objet :} Candidature en alternance -- Agents IA \& Automatisation

\vspace{0.8cm}

Madame, Monsieur,

Quatre agents IA qui raisonnent, planifient et résolvent leurs conflits pour produire des recommandations contextuelles : c'est le système que j'ai conçu avec Google ADK, LangChain et RAG, en orchestrant Gemini comme moteur d'inférence. En parallèle, mon outil open source Laravel API Generator (+30 étoiles GitHub) automatise la génération complète de code backend depuis des diagrammes UML, en intégrant un LLM pour l'assistance intelligente. Ces deux projets illustrent la convergence agents IA et automatisation du développement que Wevioo cherche à déployer.

Mon stage à INRIA Saclay (équipe TRIBE, avril à août 2026) renforce cette orientation. J'y développe en Python un pipeline d'analyse automatisée à grande échelle : NLP sur des corpus textuels (BERT/Transformers), analyse statique de SDK, et construction d'un score de confiance interprétable. Ce travail m'a habitué à structurer du code Python maintenable, à manipuler des modèles de langage en production, et à documenter rigoureusement en anglais, compétences que je mettrai directement au service du développement et de l'optimisation des agents conversationnels chez Wevioo.

Côté développement, trois ans d'expérience professionnelle en backend (Python, JavaScript, TypeScript, Docker, Git) complètent mon profil : plateforme fintech sécurisée chez CSC (Laravel, Docker, Redis), architecture microservices (NestJS, RabbitMQ, MongoDB) et plateforme de télémédecine maintenue deux ans (SOSDOCTA). Mon passage au MILA (Institut Québécois d'IA), en environnement de recherche international et anglophone, m'a appris l'autonomie et la rigueur qu'exige un contexte multiculturel comme celui de Wevioo.

Je suis disponible dès septembre 2026, à l'issue de mon stage INRIA, pour une alternance de 12 mois. Contribuer au déploiement d'agents IA qui augmentent concrètement la productivité des équipes de développement est un défi technique et humain qui me motive profondément.

Je reste à votre disposition pour un entretien afin de discuter de ma candidature.

Je vous prie d'agréer, Madame, Monsieur, l'expression de mes salutations distinguées.

\vspace{0.8cm}

\textbf{Loïc Aron MBASSI EWOLO}\\
\textit{Élève-Ingénieur ENSEA / Polytechnique Yaoundé}

\end{document}
